\subsection{Two-torsion and fibrations at the characteristic-two boundary}
\label{subsec:smallest-endpoint-geometry}

At logarithmic order two, the boundary forces two independent actual
two-torsion classes on the distinguished curve.  The same local geometry
also gives a fibration on the resolved covering surface at every odd
half-order.  The latter conclusion uses the triviality of its actual
normal line and a positive-characteristic moving-divisor criterion.

\begin{proposition}[Two-torsion at logarithmic order two]
\label{prop:smallest-endpoint-rank}
Let $k$ be algebraically closed of characteristic two, let $X$ be a
normal integral projective surface, and let $C\subset X$ be a smooth
integral effective Cartier divisor in the smooth locus, with $d=C^2>0$.
Let $M\ne\OO_X$ have order two and satisfy $M|_C\simeq\OO_C$ and
$[M]\in2\operatorname{NS}(X)$.  Suppose its logarithmic form $\omega$
has exact order two along $C$ and
\[
 2g(C)-2=3d.
\]
Put $N=\OO_C(C)$ and $\mathcal T=\omega_C\otimes N^{-3}$, and let
$Z$ be the zero divisor of the tangential leading coefficient
\[
 \alpha\in H^0(C,\omega_C\otimes N^{-2})
\]
of $\omega$.  Then $\deg Z=d$, every multiplicity of $Z$ is even,
and the logarithmic form of $\mathcal T$ satisfies
\begin{equation}\label{eq:smallest-endpoint-log-divisor}
 \operatorname{div}(\omega_{\mathcal T})=3Z.
\end{equation}
For every square root $L^{\otimes2}\simeq M$, the line bundle
$U=L|_C$ has exact order two and differs from $\mathcal T$.
Its logarithmic form $\beta_U$ is nonzero at every point of
$\operatorname{Supp}Z$.  In particular,
\[
 \dim_{\F_2}\Pic(C)[2](k)\geq2,
\]
where the group consists of actual line-bundle classes.
\end{proposition}

\begin{proof}
The lifting argument in Theorem~\ref{thm:positive-excess-picard}
supplies a square root $L$, and every such root has exact order four.
Resolve $X$ away from $C$ and retain the notation for the pullbacks.
The curve, its normal line and the logarithmic order are unchanged.

We use the geometric construction in the proof of
Corollary~\ref{cor:positive-excess-characteristic-two-boundary}.
Let $\pi:W\to X$ be the normalized $M$-torsor, let $B=\pi^*C$,
and let $D\subset S$ be its smooth strict transform on the suitable
resolution constructed there.  Write $q:D\to C$ for the induced
degree-two purely inseparable map.  Near $B$, the surface $W$ is smooth
and has the local description
\begin{equation}\label{eq:smallest-endpoint-local-model}
 \begin{aligned}
 u&=b^2+t^2A,& z&=b+ty,& y^2&=A,\\
 \omega&=t^2\frac{dA}{u},& \eta&=\frac{dz}{z}.
 \end{aligned}
\end{equation}
Here $t=0$ defines $C$, and $b,z$ are units.  The form $\eta$ is
regular on $W$ and pulls back regularly to $S$.  The calculations
\eqref{eq:boundary-conormal-order}--\eqref{eq:boundary-conormal-section}
at $a=1$ and the normal-line identification
\eqref{eq:boundary-normal-line} give
\begin{equation}\label{eq:smallest-endpoint-conormal}
 D^2=0,\qquad N_{D/S}^{-1}\simeq q^*\mathcal T,
 \qquad
 0\ne\sigma\in H^0(D,N_{D/S}^{-1}),
\end{equation}
where $\sigma$ is the conormal component of the restriction of $\eta$
as a section of $\Omega_S^1|_D$, and is nowhere zero.
These calculations apply to each choice of $L$.

The section $\alpha$ is nonzero by
Lemma~\ref{lem:positive-square-logarithmic}.  Locally its differential
representative is $d(A|_C)/b_0^2$, where $b_0=b|_C$.
Equation~\eqref{eq:boundary-leading-divisor} gives $\deg Z=d$.
The local characteristic-two derivative calculation in that proof
shows that its multiplicities are even.  Above a point $P$ of
multiplicity $2\ell>0$, the curve $B$ has complete equation
$w^2=x^{2\ell+1}$.  In particular $Z$ is nonempty and each such point
corresponds to a cusp of $B$.

\smallskip\noindent
\emph{The logarithmic form of the restricted root.}
Let $\widetilde P\in B$ lie above $P\in\operatorname{Supp}Z$.
Since $B$ is singular at $\widetilde P$ in the smooth surface $W$,
its local equation $t$ has zero cotangent vector there.
The cotangent vector $\eta(\widetilde P)$ is nonzero: otherwise its
pullback to the point of $D$ above it would vanish, contrary to
\eqref{eq:smallest-endpoint-conormal} and the injective smooth
conormal sequence on $D$.
Choose local parameters $x,t$ on $X$.  At $\widetilde P$,
\[
 dz=db+y\,dt+t\,dy=db=b_x\,dx.
\]
It follows that $b_x(P)\ne0$.  Also $A_x(P)=0$, and smoothness of
$y^2=A(x,t)$ implies $A_t(P)\ne0$.

If local frames of $L$ differ by $g_{ij}$, then
$b_i|_C=g_{ij}^2b_j|_C$.  These units give the square trivialization
of $U=L|_C$, and its logarithmic form is
\[
 \beta_U=d\log(b_i|_C).
\]
The preceding calculation proves that $\beta_U(P)\ne0$ at every
point of $\operatorname{Supp}Z$.  Thus $U$ is nontrivial, since the
logarithmic form of a trivial line bundle on a proper integral curve
is zero.  Its square is trivial, so it has exact order two.

\smallskip\noindent
\emph{The logarithmic form of the discrepancy line.}
At the generic point of $C$, write
\[
 db=R\,dt+S_1\,dA,\qquad
 R=b_t-b_x\frac{A_t}{A_x}.
\]
The relation $dA=0$ on $W$ gives
\begin{equation}\label{eq:smallest-endpoint-root-form}
 \eta=\frac{(R+y)\,dt+t\,dy}{b+ty}.
\end{equation}
Set $A_0=A|_C$ and let $R_0$ denote the rational restriction of $R$
to $C$.  If $P$ has multiplicity $2\ell$ in $Z$, the functions
$b_x|_C$ and $A_t|_C$ are units at $P$, whereas $A_x|_C$ has order
$2\ell$.  The numerator $b_tA_x-b_xA_t$ is therefore a unit, so
\[
 \operatorname{ord}_P R_0=-2\ell.
\]

We use the conormal-map identification in
\eqref{eq:boundary-normal-line}.  If $e$ is the conormal frame on
$C$ defined by $t$, the map
$q^*N^{-1}\to N_{D/S}^{-1}$ has zero divisor $q^*Z$.
Consequently the frame corresponding to $e/x^{2\ell}$ in
$N^{-1}(Z)$ is the image of $q^*e/q^*(x^{2\ell})$.
The section $\alpha$ identifies $N^{-1}(Z)$ with $\mathcal T$.
Relative to the pullback of this frame,
\eqref{eq:smallest-endpoint-root-form} gives the unit coefficient
\[
 F=q^*(x^{2\ell})\frac{R_0+y_0}{b_0},
 \qquad y_0^2=q^*A_0.
\]
Any unit introduced by the blowup chart occurs both in the conormal
map and in the chosen frame, and cancels in this coefficient.
Its square is the pullback of
\begin{equation}\label{eq:smallest-endpoint-discrepancy-unit}
 V_0=x^{4\ell}\frac{R_0^2+A_0}{b_0^2}.
\end{equation}
More generally, in any local frame of $\mathcal T$, the square of
the coefficient of $\sigma$ belongs to $q^*k(C)$, because
$k(D)^2=k(C)$.  It descends to a regular unit: its pullback is a
unit, and $q$ is finite and surjective between smooth curves.
The transition rule is that of a section of $\mathcal T^{\otimes2}$.
Thus $\sigma^2$ descends to a nowhere-zero section of this line.
The resulting square trivialization has local value $V_0^{-1}$;
in characteristic two, $d\log(V_0^{-1})=d\log V_0$.

Differentiating \eqref{eq:smallest-endpoint-discrepancy-unit} in
characteristic two gives
\begin{equation}\label{eq:smallest-endpoint-discrepancy-log}
 \omega_{\mathcal T}=d\log V_0
      =\frac{dA_0}{R_0^2+A_0}.
\end{equation}
The numerator has order $2\ell$ at $P$, and the denominator has
pole order $4\ell$, since $A_0$ is regular.  Hence
$\operatorname{ord}_P\omega_{\mathcal T}=6\ell$.
This differential is nonzero because $dA_0\ne0$ generically, and it
is regular everywhere by its construction from the torsion line.
The zeros just calculated have total degree
$3\deg Z=3d=2g(C)-2$.  They therefore exhaust its zero divisor,
proving \eqref{eq:smallest-endpoint-log-divisor}.

At every point of the nonempty divisor $Z$, the form
$\omega_{\mathcal T}$ vanishes and $\beta_U$ does not.
Any two square trivializations differ by a nonzero constant and
therefore define the same logarithmic form.  Thus $U\ne\mathcal T$.
The line $\mathcal T$ has exact order two by
Corollary~\ref{cor:positive-excess-characteristic-two-boundary}.
The two distinct nonzero classes consequently generate a subgroup
of order four, as asserted.
\end{proof}

\begin{corollary}\label{cor:smallest-endpoint-rank-strictness}
Under the surface, curve and torsion hypotheses of
Proposition~\ref{prop:smallest-endpoint-rank}, retain logarithmic
order two and $[M]\in2\operatorname{NS}(X)$, but do not assume
$2g(C)-2=3d$.  If $\dim_{\F_2}\Pic(C)[2](k)\leq1$, then
\[
 2g(C)-2-2d>d.
\]
\end{corollary}

\begin{proof}
Theorem~\ref{thm:positive-excess-picard} gives
$2g(C)-2-2d\geq d$.  Equality would contradict
Proposition~\ref{prop:smallest-endpoint-rank}.
\end{proof}

The moving-divisor input is the trivial-normal case of Totaro's
criterion \cite[Lemma~4.1, pp.~7--8]{TotaroMovingDivisors}.
The following proof also gives a multiple that is a power of the
characteristic.

\begin{lemma}[Trivial normal bundle in positive characteristic]
\label{lem:trivial-normal-moving-divisor}
Let $Y$ be a normal integral projective variety over an algebraically
closed field of characteristic $p>0$, and let $A$ be an effective
Cartier divisor with $\OO_A(A)\simeq\OO_A$.  There is an integer
$e\geq0$ such that $\OO_Y(p^eA)$ is globally generated.
\end{lemma}

\begin{proof}
Choose the indicated trivialization.  The exact sequence
\[
 0\longrightarrow\OO_Y\longrightarrow\OO_Y(A)
   \longrightarrow\OO_A\longrightarrow0
\]
gives an obstruction $c\in H^1(Y,\OO_Y)$ to lifting
$1\in H^0(A,\OO_A)$.  If $c=0$, use the identity cover below and
put $e=0$.  Otherwise, finite-dimensionality of $H^1(Y,\OO_Y)$
gives a relation among the iterates of its Frobenius-semilinear map,
\[
 \sum_{i=0}^e a_iF^ic=0,\qquad a_e=1.
\]
The highest nonzero coefficient can be normalized to one.
Thus $P(T)=\sum_{i=0}^e a_iT^{p^i}$ is a monic additive polynomial
of degree $p^e$ with $P(c)=0$.

When $c\ne0$, represent it by an additive Cech cocycle $c_{ij}$ on a
finite affine cover with affine intersections.  There are regular
functions $b_i$ such that $P(c_{ij})=b_j-b_i$.  The schemes
\[
 Y'_i=\Spec_{U_i}
       \bigl(\OO_{U_i}[T_i]/(P(T_i)-b_i)\bigr)
\]
glue by $T_i=T_j-c_{ij}$: additivity gives
$P(T_j-c_{ij})=b_j-P(c_{ij})=b_i$, and the cocycle identity gives
compatibility on triple intersections.  The monic equations define
a finite locally free faithfully flat cover
$\nu:Y'\to Y$ of degree $p^e$.  The pulled-back cocycle is a
coboundary, so $\nu^*c=0$.  In either case we have such a finite
locally free cover of degree $p^e$ killing the obstruction.

Flatness preserves the displayed exact sequence.  Hence
$\nu^*\OO_Y(A)$ has a section $s$ restricting to $1$ on $\nu^*A$.
Together with the canonical section $t$ of that effective divisor,
it has no common zero.  For each closed point $y\in Y$, choose
$\lambda\in k$ so that $s+\lambda t$ is nonzero at every point of
the finite fibre $Y'_y$.  There are only finitely many forbidden
scalars, since $s$ and $t$ have no common zero, and $k$ is infinite.
In a local trivialization, the coefficient of this section is a unit
in the Artinian fibre algebra; this remains true for a nonreduced
fibre.

The finite-flat norm, defined locally by the determinant of
multiplication by this coefficient, is therefore nonzero at $y$.
Its transition factor is the $p^e$th power of that of $\OO_Y(A)$,
so it is a global section of $\OO_Y(p^eA)$.
Such a section exists at every closed point.  The base locus is
therefore empty, proving global generation.
\end{proof}

\begin{proposition}[A fibration on the resolved cover]
\label{prop:boundary-resolved-fibration}
Under the hypotheses of
Corollary~\ref{cor:positive-excess-characteristic-two-boundary},
write $C=Y$ and $m=2a$, with $a\geq1$ odd.
Use the normalized torsor and the suitable resolution of its inverse
curve constructed in that proof, and let $D\subset S$ be the smooth
strict transform.  Then $N_{D/S}\simeq\OO_D$, and there are a
morphism with connected fibres
$g:S\to\PP^1$, a point $b\in\PP^1(k)$, and an integer $s\geq0$
such that
\begin{equation}\label{eq:boundary-fibration-fibre}
 g_*\OO_S=\OO_{\PP^1},\qquad g^*b=2^sD.
\end{equation}
This conclusion requires only $C^2>0$, not ampleness of $C$.
\end{proposition}

\begin{proof}
Equations~\eqref{eq:boundary-nonnegative-square}--\eqref{eq:boundary-conormal-section}
and their proof give $D^2=0$ and a nowhere-zero section
of $N_{D/S}^{-(h+1)}$, where $0\leq h\leq a-1$.
The cusp calculation in the same proof applies before its division
into the cases $a=1$ and $a\geq3$.  At a cusp of multiplicity
$2\ell$ in $Z$, the conormal coefficient after resolution is
\[
 v^{4\ell h}A(v^2,v^{2\ell+1}),
\]
with $A$ regular.  Since $\deg Z=d>0$, there is such a cusp;
the nowhere-zero property forces $h=0$.  The section therefore
trivializes $N_{D/S}^{-1}$ itself.

Lemma~\ref{lem:trivial-normal-moving-divisor} makes
$\OO_S(2^eD)$ globally generated for some $e\geq0$.
This line bundle is numerically nontrivial, because the nonzero
effective curve $D$ has positive intersection with an ample divisor
on $S$, and it has square zero.  Its associated morphism has
one-dimensional image: image dimension zero would make the line
trivial, whereas image dimension two would make its square positive.
Take its Stein factorization $g:S\to H$.  The base $H$ is a smooth
integral projective curve, $g_*\OO_S=\OO_H$, and
\[
 \OO_S(2^eD)\simeq g^*A_H
\]
for an ample line bundle $A_H$ on $H$.

By projection formula, the canonical section with divisor $2^eD$
descends to a section of $A_H$.  Its zero divisor is supported at one
point $b$, and the fibre over $b$ is supported on $D$.  Thus
\[
 g^*b=rD,\qquad 2^e=r\deg A_H.
\]
Both factors on the right are positive integers, so $r=2^s$ for some
$s\geq0$.

At a cusp $w^2=x^{2\ell+1}$ the final resolution chart is
\[
 w=x^\ell v,\qquad D:(x-v^2)=0.
\]
The last exceptional curve $R\simeq\PP^1$ has equation $x=0$ and
meets $D$ only at this point, with intersection length
$\operatorname{length}(k[[v]]/(v^2))=2$.
The resolution away from the inverse curve and the blowups at other
cusps leave $R$ unchanged.  The restriction of $g^*\OO_H(b)$ to $R$
has degree $2r>0$, so $R$ dominates $H$, with degree $2r$.
L\"uroth's theorem gives $H\simeq\PP^1$, also in characteristic two.
This proves \eqref{eq:boundary-fibration-fibre}.
\end{proof}

The morphism is defined on the covering resolution $S$.  Its generic
fibre may have positive genus, and its distinguished fibre may be
multiple.  These conclusions give necessary restrictions on the
boundary configuration and do not establish its existence.
